\documentclass[a4paper,11pt]{amsart}

\usepackage[T1]{fontenc}
\usepackage{lmodern}
\usepackage[margin=1in]{geometry}
\usepackage{amsmath,amssymb,mathtools}
\usepackage{microtype}
\usepackage{enumitem}
\usepackage{hyperref}
\hypersetup{
  hidelinks,
  pdftitle={Integral Solutions of z2-xy2=x3+2},
  pdfauthor={Jamal Agbanwa},
  pdfsubject={A complete effective classification of the integer solutions},
  pdfkeywords={Diophantine equation, Pell equation, generalized Pell equation}
}

\allowdisplaybreaks
\numberwithin{equation}{section}

\newtheorem{theorem}{Theorem}[section]
\newtheorem{proposition}[theorem]{Proposition}
\newtheorem{lemma}[theorem]{Lemma}
\newtheorem{corollary}[theorem]{Corollary}
\theoremstyle{definition}
\newtheorem{definition}[theorem]{Definition}
\theoremstyle{remark}
\newtheorem{remark}[theorem]{Remark}
\newtheorem{example}[theorem]{Example}

\newcommand{\Z}{\mathbb Z}
\newcommand{\R}{\mathbb R}
\newcommand{\N}{\mathbb N}
\newcommand{\eps}{\varepsilon}
\newcommand{\Acal}{\mathcal A}
\newcommand{\Rcal}{\mathcal R}
\newcommand{\Norm}{\operatorname N}

\title[Integral solutions of a cubic norm equation]
{Integral Solutions of \texorpdfstring{$z^2-xy^2=x^3+2$}{z2-xy2=x3+2}}

\author{}
\address{}
\date{}

\subjclass[2020]{Primary 11D09; Secondary 11D25, 11D41}
\keywords{Diophantine equation, Pell equation, generalized Pell equation, integral points, quadratic norm}

\begin{document}

\begin{abstract}
We give a complete and effective description of the integer solutions of
\(z^2-xy^2=x^3+2\).  The nonpositive and positive even values of \(x\) are
settled completely, while positive square values of \(x\) reduce to a finite
divisor factorization.  For every positive nonsquare value of \(x\), the
solutions split uniquely into finitely many orbits under the positive
norm-one units of \(\Z[\sqrt{x}]\), and the orbit representatives lie in an
explicit bounded interval.  In particular, a single value, \(x=23\), already
gives infinitely many integer solutions.
\end{abstract}

\maketitle

\section{Statement of the classification}

We study the equation
\begin{equation}\label{eq:main}
 z^2-xy^2=x^3+2,
 \qquad (x,y,z)\in\Z^3.
\end{equation}
The equation is unchanged by either of the independent sign changes
\(y\mapsto-y\) and \(z\mapsto-z\).  It is therefore natural to classify the
pairs \((|y|,|z|)\) and restore signs at the end.

Define
\begin{equation}\label{eq:Adef}
 \Acal=
 \left\{
 n\in\Z_{>0}:
 n\text{ is odd and every prime }p\mid n
 \text{ satisfies }p\equiv1\text{ or }7\pmod 8
 \right\}.
\end{equation}
The integer \(1\) belongs to \(\Acal\), since the condition on its prime
divisors is vacuous.

For a positive nonsquare integer \(D\), let
\begin{equation}\label{eq:fundamental-unit}
 \eps_D=u_D+v_D\sqrt D>1
\end{equation}
be the least real number of the indicated form for which
\begin{equation}\label{eq:pell-unit}
 u_D,v_D\in\Z,
 \qquad
 u_D^2-Dv_D^2=1.
\end{equation}
We shall prove from first principles that \(\eps_D\) exists, that
\(u_D,v_D>0\), and that every positive norm-one unit in \(\Z[\sqrt D]\) is a
nonnegative power of \(\eps_D\).

Put
\begin{equation}\label{eq:ND}
 N_D=D^3+2
\end{equation}
and define the reduced seed set
\begin{equation}\label{eq:seed-set}
 \Rcal_D=
 \left\{
 (A,B)\in\Z_{>0}\times\Z_{\geq0}:
 \begin{array}{l}
 A^2-DB^2=N_D,\\[2mm]
 \sqrt{N_D}\leq A+B\sqrt D<\eps_D\sqrt{N_D}
 \end{array}
 \right\}.
\end{equation}

\begin{theorem}[Complete classification]\label{thm:main}
All integer solutions of \eqref{eq:main} are exactly the following.

\begin{enumerate}[label=\textup{(\roman*)},leftmargin=2.2em]
\item The solutions with \(x\leq0\) are
\begin{equation}\label{eq:negative-solutions}
 (-1,0,1),\quad (-1,0,-1),\quad
 (-1,1,0),\quad (-1,-1,0).
\end{equation}
There are no solutions with \(x=0\) or \(x\leq-2\).

\item Let \(m\in\Acal\), put \(x=m^2\), and set
\begin{equation}\label{eq:Nm}
 N_m=m^6+2.
\end{equation}
For each positive divisor \(d\mid N_m\) satisfying
\begin{equation}\label{eq:divisor-condition}
 d<\sqrt{N_m},
 \qquad
 \frac{N_m}{d}\equiv d\pmod{2m},
\end{equation}
put
\begin{equation}\label{eq:square-formulas}
 Y_{m,d}=\frac{N_m/d-d}{2m},
 \qquad
 Z_{m,d}=\frac{N_m/d+d}{2}.
\end{equation}
Then
\begin{equation}\label{eq:square-solutions}
 (x,y,z)=\bigl(m^2,\,\sigma Y_{m,d},\,\tau Z_{m,d}\bigr),
 \qquad \sigma,\tau\in\{1,-1\},
\end{equation}
are precisely the solutions for which \(x\) is a positive square.  For fixed
\(m\), the divisor \(d\) is uniquely determined by \((|y|,|z|)\).

\item Let \(D\in\Acal\) be nonsquare.  For a seed
\((A_0,B_0)\in\Rcal_D\), define integer sequences by
\begin{equation}\label{eq:pell-recurrence}
 \begin{aligned}
 A_{n+1}&=u_DA_n+Dv_DB_n,\\
 B_{n+1}&=v_DA_n+u_DB_n
 \end{aligned}
 \qquad(n\geq0).
\end{equation}
Equivalently,
\begin{equation}\label{eq:orbit-product}
 A_n+B_n\sqrt D=(A_0+B_0\sqrt D)\eps_D^n.
\end{equation}
Then
\begin{equation}\label{eq:nonsquare-solutions}
 (x,y,z)=\bigl(D,\,\sigma B_n,\,\tau A_n\bigr),
 \qquad
 (A_0,B_0)\in\Rcal_D,
 \quad n\in\Z_{\geq0},
 \quad \sigma,\tau\in\{1,-1\},
\end{equation}
are precisely the solutions for which \(x\) is positive and nonsquare.
Every pair \((|z|,|y|)\) has a unique representation
\eqref{eq:orbit-product}.  Moreover, \(\Rcal_D\) is finite and every one of
its elements satisfies
\begin{equation}\label{eq:seed-bounds}
 0\leq B_0<v_D\sqrt{N_D},
 \qquad
 0<A_0<u_D\sqrt{N_D}.
\end{equation}
\end{enumerate}

There are no further solutions.  In particular, if \(x>0\) occurs in a
solution, then \(x\in\Acal\).
\end{theorem}

The theorem is an effective classification.  In the square case, one tests
the finitely many divisors of \(m^6+2\).  In the nonsquare case, one first
finds the least Pell unit \(\eps_D\), and then the strict bound
\(B_0^2<v_D^2N_D\) leaves only finitely many possible integers \(B_0\) to
check in \eqref{eq:seed-set}.

\section{Restrictions on the parameter \texorpdfstring{$x$}{x}}

We begin with the nonpositive values.

\begin{lemma}\label{lem:nonpositive}
The only solutions of \eqref{eq:main} with \(x\leq0\) are those in
\eqref{eq:negative-solutions}.
\end{lemma}

\begin{proof}
If \(x=0\), then \eqref{eq:main} becomes \(z^2=2\), which has no integer
solution.

Now let \(x=-a\), where \(a\in\Z_{>0}\).  Substitution in
\eqref{eq:main} gives
\begin{equation}\label{eq:negative-x}
 z^2+ay^2=2-a^3.
\end{equation}
If \(a\geq2\), then the left-hand side is nonnegative whereas
\(2-a^3\leq-6\), a contradiction.  Hence \(a=1\), so that \(x=-1\) and
\begin{equation}\label{eq:unit-circle-integral}
 y^2+z^2=1.
\end{equation}
Since \(y^2,z^2\) are nonnegative integers whose sum is \(1\), one of them
is \(1\) and the other is \(0\).  This gives exactly the four triples in
\eqref{eq:negative-solutions}.
\end{proof}

The next elementary residue lemma is the only quadratic-residue input needed
below.

\begin{lemma}\label{lem:two-residue}
Let \(p\) be an odd prime.  If the congruence
\begin{equation}\label{eq:two-square-mod-p}
 r^2\equiv2\pmod p
\end{equation}
has an integer solution, then
\begin{equation}\label{eq:pmod8}
 p\equiv1\quad\text{or}\quad7\pmod8.
\end{equation}
\end{lemma}

\begin{proof}
Write \(h=(p-1)/2\).  For each \(j\in\{1,\ldots,h\}\), reduce \(2j\)
modulo \(p\) to a signed representative
\(\delta_j a_j\), where \(\delta_j\in\{1,-1\}\) and
\(a_j\in\{1,\ldots,h\}\).  The numbers \(a_1,\ldots,a_h\) form a
permutation of \(1,\ldots,h\).  Indeed, if
\(a_i=a_j\), then \(2i\equiv\pm2j\pmod p\).  The plus sign gives \(i=j\),
whereas the minus sign would give \(i+j\equiv0\pmod p\), which is impossible
because \(2\leq i+j\leq p-1\).

Let \(M\) be the number of negative signs \(\delta_j=-1\).  Multiplying all
of the congruences \(2j\equiv\delta_ja_j\pmod p\), and then cancelling the
nonzero residue \(h!\), yields
\begin{equation}\label{eq:gauss-two}
 2^h\equiv(-1)^M\pmod p.
\end{equation}
Because \(2j\) lies between \(2\) and \(p-1\), the signed representative is
negative precisely when \(2j>p/2\).  Therefore
\begin{equation}\label{eq:Mcount}
 M=h-\left\lfloor\frac p4\right\rfloor.
\end{equation}
A direct calculation in the four odd residue classes modulo \(8\) gives
\begin{equation}\label{eq:Mparity}
 M\equiv
 \begin{cases}
 0\pmod2,&p\equiv1,7\pmod8,\\
 1\pmod2,&p\equiv3,5\pmod8.
 \end{cases}
\end{equation}

Suppose now that \eqref{eq:two-square-mod-p} holds.  Then \(p\nmid r\).
Multiplication by \(r\) permutes the nonzero residue classes modulo \(p\), so
\begin{equation*}
 r^{p-1}(p-1)!\equiv(p-1)!\pmod p.
\end{equation*}
Cancelling \((p-1)!\) proves \(r^{p-1}\equiv1\pmod p\).  Consequently
\begin{equation}\label{eq:twoh-one}
 2^h\equiv(r^2)^h=r^{p-1}\equiv1\pmod p.
\end{equation}
Equations \eqref{eq:gauss-two}--\eqref{eq:Mparity} now force \(M\) to be
even, and hence \(p\equiv1\) or \(7\pmod8\).
\end{proof}

\begin{proposition}\label{prop:positive-restrictions}
If \((x,y,z)\) satisfies \eqref{eq:main} with \(x>0\), then \(x\) is odd and
\(x\in\Acal\).  In particular, no positive even value of \(x\) occurs.
\end{proposition}

\begin{proof}
Reduction of \eqref{eq:main} modulo \(x\) gives
\begin{equation}\label{eq:z2modx}
 z^2\equiv2\pmod x.
\end{equation}
Thus, for every odd prime \(p\mid x\), the congruence \(z^2\equiv2\pmod p\)
holds.  Lemma~\ref{lem:two-residue} then gives
\(p\equiv1\) or \(7\pmod8\).

It remains to exclude even \(x\).  If \(4\mid x\), then
\eqref{eq:z2modx} implies \(z^2\equiv2\pmod4\), impossible because a square
modulo \(4\) is \(0\) or \(1\).  Hence any hypothetical even solution would
have
\begin{equation}\label{eq:x2m}
 x=2m
\end{equation}
with \(m\) odd.  Equation \eqref{eq:main} becomes
\begin{equation}\label{eq:even-x-equation}
 z^2=2my^2+8m^3+2.
\end{equation}
If \(y\) were even, then the right-hand side would be congruent to \(2\)
modulo \(8\), again impossible for a square.  Thus \(y\) is odd.  Every odd
square is congruent to \(1\) or \(9\pmod{16}\), and therefore
\(2my^2\equiv2m\pmod{16}\).  Since \(m\) is odd, also
\(8m^3\equiv8\pmod{16}\).  It follows that
\begin{equation}\label{eq:z2mod16}
 z^2\equiv2m+10\pmod{16}.
\end{equation}
The right-hand side is even, so it must be one of the even square residues
\(0,4\pmod{16}\).  Hence
\begin{equation}\label{eq:m35}
 m\equiv3\quad\text{or}\quad5\pmod8.
\end{equation}

On the other hand, reducing \eqref{eq:even-x-equation} modulo \(m\) gives
\(z^2\equiv2\pmod m\).  By the first paragraph of the proof, every prime
divisor of \(m\) is congruent to \(1\) or \(-1\pmod8\).  The product of
arbitrarily many residues from \(\{1,-1\}\) is again \(1\) or \(-1\), even
when prime powers are counted with multiplicity.  Therefore
\begin{equation}\label{eq:mpm1}
 m\equiv1\quad\text{or}\quad7\pmod8,
\end{equation}
contradicting \eqref{eq:m35}.  Thus \(x\) is odd.  Together with the prime
restriction already proved, this says exactly that \(x\in\Acal\).
\end{proof}

\section{The positive square case}

We now give a finite factorization of all solutions for which \(x\) is a
square.

\begin{proposition}\label{prop:square-case}
Let \(m\in\Z_{>0}\).  Equation \eqref{eq:main} has a solution with
\(x=m^2\) if and only if \(m\in\Acal\) and there is a positive divisor
\(d\mid m^6+2\) satisfying \eqref{eq:divisor-condition}.  For every such
\(d\), all corresponding solutions are given by
\eqref{eq:square-formulas}--\eqref{eq:square-solutions}, and this
parametrization is unique up to the independent signs of \(y\) and \(z\).
\end{proposition}

\begin{proof}
Suppose first that \((m^2,y,z)\) is a solution.  By
Proposition~\ref{prop:positive-restrictions}, \(m^2\in\Acal\), and hence
\(m\in\Acal\).  In particular, \(m\) is odd.  The equation is
\begin{equation}\label{eq:square-norm}
 z^2-m^2y^2=m^6+2=N_m.
\end{equation}
The integer \(N_m\) is not a square, since
\begin{equation}\label{eq:Nm-between-squares}
 m^6<N_m=m^6+2<(m^3+1)^2=m^6+2m^3+1.
\end{equation}
In particular, \(y\neq0\), for otherwise \eqref{eq:square-norm} would make
\(N_m=z^2\).

Set
\begin{equation}\label{eq:YZabs}
 Y=|y|,
 \qquad Z=|z|.
\end{equation}
Since \(Z^2-m^2Y^2=N_m>0\), we have \(Z>mY\).  Therefore
\begin{equation}\label{eq:de-def}
 d=Z-mY,
 \qquad
 e=Z+mY
\end{equation}
are positive integers satisfying
\begin{equation}\label{eq:de-properties}
 de=N_m,
 \qquad d<e,
 \qquad e-d=2mY.
\end{equation}
Thus \(d<\sqrt{N_m}\), \(e=N_m/d\), and
\begin{equation}\label{eq:econgd}
 \frac{N_m}{d}=e\equiv d\pmod{2m}.
\end{equation}
Solving \eqref{eq:de-def} for \(Y,Z\) gives
\begin{equation}\label{eq:YZ-from-de}
 Y=\frac{e-d}{2m}=\frac{N_m/d-d}{2m},
 \qquad
 Z=\frac{e+d}{2}=\frac{N_m/d+d}{2}.
\end{equation}
This proves the necessity of the stated formula.  It also proves uniqueness:
from \((Y,Z)\), the divisor is forced to be \(d=Z-mY\).

Conversely, suppose \(m\in\Acal\) and \(d\mid N_m\) satisfies
\eqref{eq:divisor-condition}.  Put \(e=N_m/d\).  Since \(m\) is odd,
\(N_m\), \(d\), and \(e\) are all odd.  The congruence
\(e\equiv d\pmod{2m}\) shows that
\begin{equation*}
 Y=\frac{e-d}{2m}
\end{equation*}
is an integer, and the oddness of \(d,e\) shows that
\begin{equation*}
 Z=\frac{e+d}{2}
\end{equation*}
is an integer.  The strict inequality \(d<\sqrt{N_m}\) implies \(e>d\), so
\(Y>0\).  Finally,
\begin{align}\label{eq:factor-verification}
 Z^2-m^2Y^2
 &=\frac{(e+d)^2-(e-d)^2}{4}\notag\\
 &=ed=N_m=m^6+2.
\end{align}
This is precisely \eqref{eq:main} with \(x=m^2\).  Independent choices of
the signs of \(y\) and \(z\) give all solutions belonging to \((Y,Z)\).
\end{proof}

\begin{example}\label{ex:square}
For \(m=1\), one has \(N_m=3\), and \(d=1\) gives
\begin{equation*}
 (x,y,z)=(1,\pm1,\pm2).
\end{equation*}
For \(m=7\), one has \(N_m=7^6+2=117651\).  Taking \(d=3\), so that
\(N_m/d=39217\), gives
\begin{equation*}
 Y=\frac{39217-3}{14}=2801,
 \qquad
 Z=\frac{39217+3}{2}=19610.
\end{equation*}
Hence
\begin{equation*}
 (x,y,z)=(49,\pm2801,\pm19610)
\end{equation*}
is a fourfold sign family of solutions.
\end{example}

\section{Pell units and the nonsquare case}

For a positive nonsquare \(D\), write
\begin{equation}\label{eq:norm-definition}
 \Norm(a+b\sqrt D)=(a+b\sqrt D)(a-b\sqrt D)=a^2-Db^2.
\end{equation}
We first prove all of the unit theory that will be used.

\begin{proposition}[Positive Pell units]\label{prop:pell-units}
Let \(D\in\Z_{>0}\) be nonsquare.  Then the following statements hold.

\begin{enumerate}[label=\textup{(\alph*)},leftmargin=2em]
\item There exist positive integers \(u,v\) such that
\begin{equation}\label{eq:pell-existence}
 u^2-Dv^2=1.
\end{equation}

\item Among all real numbers \(u+v\sqrt D>1\) with \(u,v\in\Z\) and
\(u^2-Dv^2=1\), there is a least one, denoted
\(\eps_D=u_D+v_D\sqrt D\).  Its coefficients satisfy
\(u_D,v_D\in\Z_{>0}\).

\item If \(\eta=a+b\sqrt D>1\), where \(a,b\in\Z\) and
\(a^2-Db^2=1\), then there is a unique integer \(n\geq1\) such that
\begin{equation}\label{eq:unit-power}
 \eta=\eps_D^n.
\end{equation}
Together with \(1=\eps_D^0\), the nonnegative powers of \(\eps_D\) are
exactly the positive norm-one units whose real value is at least \(1\).
\end{enumerate}
\end{proposition}

\begin{proof}
We first prove existence.  For every integer \(T\geq2\), consider the
\(T+1\) fractional parts
\begin{equation*}
 0,\ \{\sqrt D\},\ \{2\sqrt D\},\ldots,\ \{T\sqrt D\}.
\end{equation*}
Divide \([0,1)\) into \(T\) half-open intervals of length \(1/T\).  Two of
the displayed fractional parts lie in the same interval.  Subtracting the
corresponding multiples of \(\sqrt D\), we obtain integers \(r,s\) such that
\begin{equation}\label{eq:dirichlet-approx}
 1\leq s\leq T,
 \qquad
 |s\sqrt D-r|<\frac1T.
\end{equation}
Since \(s\sqrt D\geq\sqrt2\) and \(1/T\leq1/2\), every \(r\)
satisfying \eqref{eq:dirichlet-approx} is positive.  The integer
\begin{equation}\label{eq:kdef}
 k=r^2-Ds^2
\end{equation}
is nonzero because \(\sqrt D\) is irrational.  Moreover,
\begin{align}\label{eq:kbound}
 |k|
 &=|r-s\sqrt D|\,(r+s\sqrt D)\notag\\
 &<\frac1T\left(2s\sqrt D+\frac1T\right)\notag\\
 &\leq2\sqrt D+\frac1{T^2}
 <2\sqrt D+1.
\end{align}

As \(T\) tends to infinity, the denominators \(s\) obtained in
\eqref{eq:dirichlet-approx} cannot remain bounded.  Indeed, if
\(1\leq s\leq S\), then each distance
\begin{equation*}
 \operatorname{dist}(s\sqrt D,\Z)
\end{equation*}
is positive; the minimum of these finitely many positive numbers would be a
positive constant \(\delta\), contradicting
\eqref{eq:dirichlet-approx} for \(T>1/\delta\).  Hence we obtain infinitely
many distinct pairs \((r,s)\).

By \eqref{eq:kbound}, the nonzero integer \(k=r^2-Ds^2\) lies in a fixed
finite set.  Thus one fixed nonzero value of \(k\) occurs for infinitely many
pairs.  Within this infinite fixed-\(k\) subsequence, \(s\) is again unbounded,
because for each fixed \(s\) the equation \(r^2-Ds^2=k\) permits at most two
integers \(r\).  Hence \(r+s\sqrt D\) is unbounded.  Among those pairs, two,
say \((r_1,s_1)\) and \((r_2,s_2)\), have simultaneously
\begin{equation}\label{eq:residue-pairs}
 r_1\equiv r_2\pmod{|k|},
 \qquad
 s_1\equiv s_2\pmod{|k|}.
\end{equation}
We may label the two pairs so that
\begin{equation}\label{eq:alpha-order}
 r_1+s_1\sqrt D>r_2+s_2\sqrt D>0.
\end{equation}
Consider their quotient.  Since both numerator and denominator have norm
\(k\),
\begin{align}\label{eq:unit-quotient}
 \frac{r_1+s_1\sqrt D}{r_2+s_2\sqrt D}
 &=\frac{(r_1+s_1\sqrt D)(r_2-s_2\sqrt D)}{k}\notag\\
 &=u+v\sqrt D,
\end{align}
where
\begin{equation}\label{eq:uvquotient}
 u=\frac{r_1r_2-Ds_1s_2}{k},
 \qquad
 v=\frac{s_1r_2-r_1s_2}{k}.
\end{equation}
The congruences \eqref{eq:residue-pairs} imply that both numerators in
\eqref{eq:uvquotient} are divisible by \(|k|\), and hence \(u,v\in\Z\).
Taking norms in \eqref{eq:unit-quotient} gives
\begin{equation*}
 u^2-Dv^2=1.
\end{equation*}
The quotient in \eqref{eq:unit-quotient} is greater than \(1\) by
\eqref{eq:alpha-order}.  Its conjugate is its positive reciprocal.  Therefore
\begin{equation*}
 u=\frac{\eta+\eta^{-1}}2>0,
 \qquad
 v=\frac{\eta-\eta^{-1}}{2\sqrt D}>0,
\end{equation*}
where \(\eta\) denotes the quotient.  This proves part~\textup{(a)}.

Let \(\mathcal U_D\) be the nonempty set of all integer-coefficient
norm-one units with real value greater than \(1\).  Choose any
\(\eta_0\in\mathcal U_D\).  There are only finitely many elements
\(u+v\sqrt D\in\mathcal U_D\) satisfying
\(1<u+v\sqrt D\leq\eta_0\), because then \(0<v<\eta_0/\sqrt D\) and
\(0<u<\eta_0\).  That finite nonempty set has a least element; it is also the
least element of all of \(\mathcal U_D\).  Denote it by \(\eps_D\).  The
same formulas involving a unit and its positive reciprocal show that its
coefficients \(u_D,v_D\) are positive.  This proves part~\textup{(b)}.

Finally, let \(\eta\in\mathcal U_D\).  Since \(\eps_D>1\), there is a unique
integer \(n\geq0\) such that
\begin{equation}\label{eq:archimedean-reduction-unit}
 \eps_D^n\leq\eta<\eps_D^{n+1}.
\end{equation}
The inverse \(\eps_D^{-1}=u_D-v_D\sqrt D\) again has integer coefficients
and norm \(1\).  Hence
\begin{equation*}
 \theta=\eta\eps_D^{-n}
\end{equation*}
is an integer-coefficient norm-one unit satisfying
\begin{equation}\label{eq:theta-range}
 1\leq\theta<\eps_D.
\end{equation}
If \(\theta>1\), then \(\theta\in\mathcal U_D\), contradicting the
minimality of \(\eps_D\).  Thus \(\theta=1\), and
\(\eta=\eps_D^n\).  Since \(\eta>1\), one has \(n\geq1\).  Uniqueness of
\(n\) follows from the strict monotonicity of the powers of \(\eps_D>1\).
This proves part~\textup{(c)}.
\end{proof}

We next prove a general orbit decomposition for a positive generalized Pell
equation.  This is the exact mechanism behind part~\textup{(iii)} of
Theorem~\ref{thm:main}.

\begin{proposition}[Finite reduced-orbit decomposition]\label{prop:generalized-pell}
Let \(D\in\Z_{>0}\) be nonsquare, let \(N\in\Z_{>0}\), and let
\(\eps_D=u_D+v_D\sqrt D\) be as in
Proposition~\ref{prop:pell-units}.  Define
\begin{equation}\label{eq:RDN}
 \Rcal(D,N)=
 \left\{
 (A,B)\in\Z_{>0}\times\Z_{\geq0}:
 \begin{array}{l}
 A^2-DB^2=N,\\[1mm]
 \sqrt N\leq A+B\sqrt D<\eps_D\sqrt N
 \end{array}
 \right\}.
\end{equation}
Then multiplication induces a bijection
\begin{equation}\label{eq:orbit-bijection}
 \Rcal(D,N)\times\Z_{\geq0}
 \longrightarrow
 \left\{(A,B)\in\Z_{>0}\times\Z_{\geq0}:A^2-DB^2=N\right\},
\end{equation}
where \(((A_0,B_0),n)\) is sent to the coefficients of
\begin{equation}\label{eq:general-orbit}
 (A_0+B_0\sqrt D)\eps_D^n.
\end{equation}
The seed set \(\Rcal(D,N)\) is finite, and every seed satisfies
\begin{equation}\label{eq:general-seed-bounds}
 0\leq B<v_D\sqrt N,
 \qquad
 0<A<u_D\sqrt N.
\end{equation}
\end{proposition}

\begin{proof}
Let \((A,B)\in\Z_{>0}\times\Z_{\geq0}\) satisfy
\begin{equation}\label{eq:genpell-solution}
 A^2-DB^2=N.
\end{equation}
Set
\begin{equation}\label{eq:tdef}
 t=A+B\sqrt D.
\end{equation}
Since \(A>0\), \(B\geq0\), and \eqref{eq:genpell-solution} has positive
right-hand side, one has \(A>B\sqrt D\).  Thus both \(t\) and its conjugate
\(t'=A-B\sqrt D\) are positive, and
\begin{equation}\label{eq:ttprime}
 tt'=N.
\end{equation}
Moreover, \(t\geq t'\), so \(t^2\geq tt'=N\), whence
\begin{equation}\label{eq:t-lower}
 t\geq\sqrt N.
\end{equation}

There is a unique integer \(n\geq0\) for which
\begin{equation}\label{eq:nchoice}
 \eps_D^n\leq\frac{t}{\sqrt N}<\eps_D^{n+1}.
\end{equation}
Put
\begin{equation}\label{eq:t0def}
 t_0=t\eps_D^{-n}.
\end{equation}
Since \(\eps_D^{-1}=u_D-v_D\sqrt D\), the number \(t_0\) has integer
coefficients:
\begin{equation}\label{eq:t0coeff}
 t_0=A_0+B_0\sqrt D,
 \qquad A_0,B_0\in\Z.
\end{equation}
Its norm is \(N\), and \eqref{eq:nchoice} gives
\begin{equation}\label{eq:t0interval}
 \sqrt N\leq t_0<\eps_D\sqrt N.
\end{equation}
Its conjugate is \(t_0'=N/t_0>0\).  In addition,
\(t_0\geq\sqrt N\geq t_0'\), so
\begin{equation}\label{eq:A0B0positive}
 A_0=\frac{t_0+t_0'}2>0,
 \qquad
 B_0=\frac{t_0-t_0'}{2\sqrt D}\geq0.
\end{equation}
Thus \((A_0,B_0)\in\Rcal(D,N)\), and
\begin{equation*}
 t=(A_0+B_0\sqrt D)\eps_D^n.
\end{equation*}
This proves surjectivity of \eqref{eq:orbit-bijection}.

The integer \(n\) is uniquely determined by \eqref{eq:nchoice}, after which
\(t_0=t\eps_D^{-n}\) is uniquely determined.  Hence the representation is
unique, proving injectivity.

Conversely, start with a seed \((A_0,B_0)\in\Rcal(D,N)\).  Because
\(\eps_D^n\in\Z[\sqrt D]\) and has norm \(1\), the product
\eqref{eq:general-orbit} has integer coefficients \(A_n,B_n\) and norm
\(N\).  Since \(A_0>0\), \(B_0\geq0\), and \(u_D,v_D>0\), its coefficients
remain in \(\Z_{>0}\times\Z_{\geq0}\).  Explicitly they obey
\begin{equation}\label{eq:general-recurrence}
 A_{n+1}=u_DA_n+Dv_DB_n,
 \qquad
 B_{n+1}=v_DA_n+u_DB_n.
\end{equation}
This verifies directly that the map in \eqref{eq:orbit-bijection} is
well-defined.

It remains to establish the bounds and finiteness.  For a seed, let
\(t=A+B\sqrt D\).  Since its conjugate equals \(N/t\),
\begin{equation}\label{eq:ABfromt}
 A=\frac{t+N/t}{2},
 \qquad
 B=\frac{t-N/t}{2\sqrt D}.
\end{equation}
On \([\sqrt N,\infty)\), the first right-hand side is nondecreasing and the
second is strictly increasing.  Since \(t<\eps_D\sqrt N\), we obtain
\begin{align}\label{eq:A-bound-calc}
 A
 &<\frac{\eps_D+\eps_D^{-1}}2\sqrt N
 =u_D\sqrt N,
\end{align}
and
\begin{align}\label{eq:B-bound-calc}
 B
 &<\frac{\eps_D-\eps_D^{-1}}{2\sqrt D}\sqrt N
 =v_D\sqrt N.
\end{align}
Here we used
\begin{equation*}
 \eps_D^{-1}=u_D-v_D\sqrt D.
\end{equation*}
The bounds leave only finitely many integer pairs \((A,B)\), so
\(\Rcal(D,N)\) is finite.
\end{proof}

\begin{proposition}\label{prop:nonsquare-case}
Let \(D>0\) be nonsquare.  The solutions of \eqref{eq:main} with \(x=D\)
are exactly those in \eqref{eq:nonsquare-solutions}.  If there is at least
one such solution, then there are infinitely many.
\end{proposition}

\begin{proof}
For \(x=D\), equation \eqref{eq:main} is
\begin{equation}\label{eq:nonsquare-norm}
 z^2-Dy^2=D^3+2=N_D.
\end{equation}
Given a solution, put \(A=|z|\) and \(B=|y|\).  Then
\((A,B)\in\Z_{>0}\times\Z_{\geq0}\) and
\begin{equation*}
 A^2-DB^2=N_D.
\end{equation*}
Proposition~\ref{prop:generalized-pell}, applied with \(N=N_D\), gives a
unique seed \((A_0,B_0)\in\Rcal_D\) and a unique \(n\geq0\) such that
\begin{equation*}
 A+B\sqrt D=(A_0+B_0\sqrt D)\eps_D^n.
\end{equation*}
This is exactly the representation in \eqref{eq:nonsquare-solutions} after
the independent signs of \(y,z\) are restored.

Conversely, each product in \eqref{eq:orbit-product} has norm \(N_D\), so
its coefficients satisfy \eqref{eq:nonsquare-norm}; hence every listed triple
is a solution.  The bounds \eqref{eq:seed-bounds} are the specialization of
\eqref{eq:general-seed-bounds}.

If a seed exists, then the positive real numbers
\begin{equation*}
 (A_0+B_0\sqrt D)\eps_D^n
 \qquad(n\geq0)
\end{equation*}
are strictly increasing because \(\eps_D>1\).  Their coefficient pairs are
therefore distinct, and so the seed produces infinitely many solutions.
\end{proof}

\begin{proof}[Proof of Theorem~\ref{thm:main}]
Lemma~\ref{lem:nonpositive} proves part~\textup{(i)}.
Proposition~\ref{prop:positive-restrictions} shows that every positive
solution value of \(x\) belongs to \(\Acal\) and that no positive even value
is possible.  Every positive integer is either a square or a nonsquare.
Proposition~\ref{prop:square-case} gives the complete and unique description
in the square case, while Proposition~\ref{prop:nonsquare-case} and
Proposition~\ref{prop:generalized-pell} give the complete, unique, and finite
seed-orbit description in the nonsquare case.  These mutually exclusive
cases exhaust all integers \(x\), so no further solutions exist.
\end{proof}

\section{Consequences and examples}

\begin{corollary}\label{cor:fixed-x}
For each fixed integer \(x\), the following alternatives hold.
\begin{enumerate}[label=\textup{(\roman*)},leftmargin=2.2em]
\item If \(x<-1\), \(x=0\), or \(x>0\) is even, there are no solutions.
\item If \(x=-1\), there are exactly four solutions.
\item If \(x>0\) is a square, there are finitely many solutions, explicitly
given by the divisor conditions \eqref{eq:divisor-condition}.
\item If \(x>0\) is nonsquare, then there are either no solutions or
infinitely many.  In the latter case the solutions are a disjoint finite
union of Pell orbits, one orbit for each seed in \(\Rcal_x\).
\end{enumerate}
\end{corollary}

\begin{proof}
Only the last assertion requires comment.  The seed set is finite by
\eqref{eq:seed-bounds}; uniqueness in
Proposition~\ref{prop:generalized-pell} makes the corresponding orbits
disjoint; and Proposition~\ref{prop:nonsquare-case} shows that every nonempty
orbit is infinite.
\end{proof}

\begin{corollary}[An explicit infinite family]\label{cor:x23}
Define positive integers \(Y_n,Z_n\) by
\begin{equation}\label{eq:x23initial}
 (Y_0,Z_0)=(23,156)
\end{equation}
and
\begin{equation}\label{eq:x23recurrence}
 \begin{aligned}
 Z_{n+1}&=24Z_n+115Y_n,\\
 Y_{n+1}&=5Z_n+24Y_n
 \end{aligned}
 \qquad(n\geq0).
\end{equation}
Then, for every \(n\geq0\) and arbitrary signs
\(\sigma,\tau\in\{1,-1\}\),
\begin{equation}\label{eq:x23solutions}
 (x,y,z)=(23,\sigma Y_n,\tau Z_n)
\end{equation}
is a solution of \eqref{eq:main}.  These give infinitely many distinct
integer solutions.
\end{corollary}

\begin{proof}
The two identities
\begin{equation}\label{eq:x23norms}
 24^2-23\cdot5^2=1
\end{equation}
and
\begin{equation}\label{eq:x23base}
 156^2-23\cdot23^2=12169=23^3+2
\end{equation}
show that
\begin{equation}\label{eq:x23product}
 Z_n+Y_n\sqrt{23}
 =(156+23\sqrt{23})(24+5\sqrt{23})^n
\end{equation}
has norm \(23^3+2\) for every \(n\geq0\).  Equating coefficients after one
further multiplication by \(24+5\sqrt{23}\) gives exactly the recurrence
\eqref{eq:x23recurrence}.  Hence
\begin{equation*}
 Z_n^2-23Y_n^2=23^3+2,
\end{equation*}
which is \eqref{eq:main} with \(x=23\).  Moreover,
\begin{equation*}
 Y_{n+1}=5Z_n+24Y_n>Y_n,
\end{equation*}
so the pairs \((Y_n,Z_n)\) are distinct.  The first two are
\begin{equation*}
 (Y_0,Z_0)=(23,156),
 \qquad
 (Y_1,Z_1)=(1332,6389).
\end{equation*}
\end{proof}

\begin{remark}[The local condition is not sufficient]\label{rem:x7}
Membership in \(\Acal\) is a necessary restriction on a positive value of
\(x\), but it is not sufficient.  For example, \(7\in\Acal\), yet \(x=7\)
gives no solution.  Indeed, a solution would satisfy
\begin{equation}\label{eq:x7}
 z^2-7y^2=7^3+2=345.
\end{equation}
Modulo \(5\), this becomes \(z^2\equiv2y^2\pmod5\).  If \(5\nmid y\), then
\((zy^{-1})^2\equiv2\pmod5\), impossible because the square residues modulo
\(5\) are \(0,1,4\).  Thus \(5\mid y\), and then \(5\mid z\) as well.
Consequently \(25\) divides the left-hand side of \eqref{eq:x7}, whereas
\(25\nmid345\), a contradiction.  The finite seed condition
\(\Rcal_D\neq\varnothing\) is the exact remaining global condition in the
nonsquare case.
\end{remark}

\end{document}
